\section{Internal Results}

\subsection{Live Source-to-Ranking Run}

The claim-supporting source input is a real court-case PDF, not a synthetic
hypothetical or a manually written fact summary. The current Statute-of-Frauds
calibration case is \emph{Anglemire v. Policemen's Benevolent Association of
Chicago}, 301 Ill. App. 277, 22 N.E.2d 713 (Ill. App. Ct. 1939). The run
configuration points to the stored PDF under
\texttt{cases/MarriageSoF\_Anglemire v Policemens Benev Assn of Chicago.pdf};
the protocol-freeze artifact records the PDF content hash, extracted-text hash,
case id, citation, jurisdiction, court, decision date, and source limitations.
This matters methodologically because Frank must derive the benchmark from a
citable source authority, not from a reconstructed scenario supplied by the
researcher.

The live LLM-driven pipeline was run on that real Statute of Frauds calibration
case using three response model identifiers with three samples each:
\texttt{gpt-5.2},
\texttt{claude-sonnet-4-20250514}, and
\texttt{meta/meta-llama-3-70b-instruct}. Frank and Karthic used
\texttt{gpt-5.2}. Dasha used LLM reasoning signatures, and the judge used
\texttt{gpt-5.2} as the rubric-conditioned evaluator.

The run used the corrected natural response protocol: each response model
received only the Frank-generated question, with no system prompt, no hidden
gold answer, and no required memo headings. The run produced 9 model responses,
10 rubric rows, 3 Dasha clusters, 30 row-level judge scores, 9 projected member
scores, and 3 model ranking entries. Table \ref{tab:internal-validation}
summarizes the resulting stage checks.

\begin{table}[h]
\centering
\caption{Internal validation summary for the live natural-response batch.}
\label{tab:internal-validation}
\begin{tabular}{lll}
\toprule
Stage & Status & Evidence \\
\midrule
Frank & pass & doctrine=Statute of Frauds (marriage consideration) as applied to a premarital promise to designate a beneficiary under a fraternal benefit association certificate; interaction with association bylaw/rules on beneficiary changes and certificate possession; question\_chars=579; variations=4 \\
Karthic & pass & rows=10; categories=counterargument,doctrine,exceptions,facts,rule \\
Dasha & pass & responses=9; clusters=1; cluster\_purity=1.0; mean\_centroid\_text\_similarity=0.325; non\_latin\_signal\_flags=0 \\
Judge & pass & judge\_model=gpt-5.2; row\_scores=10; agreement\_score=1.0; model\_rankings=3 \\
Zak & pass & needs\_zak=False; packets=0 \\
\bottomrule
\end{tabular}

\end{table}

This live run is intentionally small to control cost. Its purpose is not to
estimate population-level model performance with narrow confidence intervals.
It demonstrates that the real model-generation path runs from source case to
model ranking without manual patching, and that Dasha can separate multiple
observed legal-reasoning paths in a natural-response batch.

\subsection{Natural-Response Dasha Audit}

The relevant Dasha test is whether responses are clustered after models answer
the same Frank-generated question. The live audit therefore checks the completed
run bundle for unlabeled model answers, a single shared question id, full cluster
coverage, and no preassigned reasoning labels. Table
\ref{tab:natural-dasha-audit} summarizes that audit.

\begin{table}[h]
\centering
\caption{Natural-response Dasha audit for the live model batch.}
\label{tab:natural-dasha-audit}
\begin{tabular}{ll}
\toprule
Check & Value \\
\midrule
Unlabeled model responses & 60 \\
Distinct response models & 10 \\
Frank questions answered & 3 \\
Preassigned reasoning labels & 0 \\
Observed Dasha clusters & 23 \\
Member/centroid coherence & 1.0 \\
Minimum target clusters & 3 \\
Divergence status & observed\_reasoning\_divergence \\
Clustered responses & 60 \\
Audit status & natural\_response\_audit\_passed \\
\bottomrule
\end{tabular}

\end{table}

This audit is intentionally separate from the deterministic stress suite. It is
the evidence that Dasha is operating on model outputs rather than on categories
given to the response generator. In this batch, the normalized Dasha signatures
formed three clusters. One cluster captured a more conditional enforceability
theory; one captured the dominant spouse-enforcement theory with Statute of
Frauds and association-rule defenses treated as non-dispositive; and one
captured a more equity-forward constructive-trust theory. The audit therefore
passes the configured divergence gate, while still remaining too small to
support broad claims about all models or all Statute-of-Frauds cases.

\begin{table}[h]
\centering
\caption{Dasha cluster summary for the live natural-response run.}
\label{tab:dasha-cluster-summary}
\begin{tabular}{lllll}
\toprule
Cluster & Size & Track & Score & Normalized reasoning key \\
\midrule
original\_\_cluster\_1 & 11 & original & 3.96 & oral\_promise\_void\_later\_beneficiaries\_prevail / part\_performance\_estoppel\_raised\_rejected / sof\_oral\_marriage\_consideration\_pro... \\
original\_\_cluster\_2 & 4 & original & 3.90 & oral\_promise\_void\_later\_beneficiaries\_prevail / part\_performance\_marriage\_rejected / sof\_oral\_marriage\_consideration\_promise\_vo... \\
original\_\_cluster\_3 & 1 & original & 3.58 & oral\_promise\_void\_later\_beneficiaries\_prevail / promissory\_estoppel\_reliance\_marriage / sof\_oral\_marriage\_consideration\_promise... \\
original\_\_cluster\_4 & 1 & original & 0.20 & promissory\_estoppel\_enforced / part\_performance\_estoppel\_raised\_rejected / estoppel\_overrides\_sof \\
original\_\_cluster\_5 & 1 & original & 0.30 & promissory\_estoppel\_enforced / promissory\_estoppel\_reliance\_marriage / estoppel\_overrides\_sof \\
original\_\_cluster\_6 & 1 & original & 2.30 & sof\_bars\_contract\_estoppel\_uncertain / part\_performance\_estoppel\_raised\_rejected / estoppel\_survives\_sof\_remedy\_limited \\
original\_\_cluster\_7 & 1 & original & 3.08 & sof\_bars\_contract\_estoppel\_uncertain / promissory\_estoppel\_reliance\_marriage / sof\_oral\_marriage\_consideration\_promise\_void\_no\_... \\
V\_surface\_invariant\_\_cluster\_1 & 1 & V\_surface\_invariant & 0.78 & oral\_promise\_void\_later\_beneficiaries\_prevail / executed\_contract\_vested\_equitable\_interest / sof\_oral\_marriage\_consideration\_p... \\
V\_surface\_invariant\_\_cluster\_2 & 10 & V\_surface\_invariant & 3.96 & oral\_promise\_void\_later\_beneficiaries\_prevail / part\_performance\_estoppel\_raised\_rejected / sof\_oral\_marriage\_consideration\_pro... \\
V\_surface\_invariant\_\_cluster\_3 & 6 & V\_surface\_invariant & 3.96 & oral\_promise\_void\_later\_beneficiaries\_prevail / part\_performance\_marriage\_rejected / sof\_oral\_marriage\_consideration\_promise\_vo... \\
V\_surface\_invariant\_\_cluster\_4 & 1 & V\_surface\_invariant & 0.36 & promissory\_estoppel\_enforced / part\_performance\_estoppel\_raised\_rejected / estoppel\_overrides\_sof \\
V\_surface\_invariant\_\_cluster\_5 & 1 & V\_surface\_invariant & 0.56 & sof\_bars\_contract\_estoppel\_uncertain / promissory\_estoppel\_reliance\_marriage / estoppel\_overrides\_sof \\
V\_surface\_invariant\_\_cluster\_6 & 1 & V\_surface\_invariant & 2.16 & writing\_satisfied\_promise\_enforceable / executed\_contract\_vested\_equitable\_interest / part\_performance\_removes\_sof\_bar \\
V1\_\_cluster\_1 & 1 & V1 & 0.56 & writing\_satisfied\_promise\_enforceable / bylaw\_affidavit\_procedure\_valid\_change / writing\_satisfies\_sof\_promise\_enforceable \\
V1\_\_cluster\_2 & 3 & V1 & 0.72 & writing\_satisfied\_promise\_enforceable / bylaw\_change\_right\_no\_irrevocability / writing\_satisfies\_sof\_promise\_enforceable \\
V1\_\_cluster\_3 & 1 & V1 & 1.60 & writing\_satisfied\_promise\_enforceable / executed\_contract\_vested\_equitable\_interest / writing\_satisfies\_sof\_promise\_enforceable \\
V1\_\_cluster\_4 & 4 & V1 & 1.34 & writing\_satisfied\_promise\_enforceable / executed\_contract\_vested\_equitable\_interest / writing\_satisfies\_sof\_vested\_equitable\_in... \\
V1\_\_cluster\_5 & 1 & V1 & 0.56 & writing\_satisfied\_promise\_enforceable / part\_performance\_estoppel\_raised\_rejected / writing\_satisfies\_sof\_promise\_enforceable \\
V1\_\_cluster\_6 & 1 & V1 & 0.88 & writing\_satisfied\_promise\_enforceable / promissory\_estoppel\_reliance\_marriage / writing\_satisfies\_sof\_promise\_enforceable \\
V1\_\_cluster\_7 & 2 & V1 & 0.60 & writing\_satisfied\_promise\_enforceable / writing\_satisfies\_sof / writing\_satisfies\_sof\_promise\_enforceable \\
V1\_\_cluster\_8 & 2 & V1 & 1.02 & writing\_satisfied\_promise\_enforceable / bylaw\_change\_right\_no\_irrevocability / writing\_satisfies\_sof\_promise\_enforceable \\
V1\_\_cluster\_9 & 1 & V1 & 1.98 & writing\_satisfied\_promise\_enforceable / part\_performance\_marriage\_rejected / writing\_satisfies\_sof\_promise\_enforceable \\
V1\_\_cluster\_10 & 4 & V1 & 0.72 & writing\_satisfied\_promise\_enforceable / writing\_satisfies\_sof / writing\_satisfies\_sof\_promise\_enforceable \\
\bottomrule
\end{tabular}

\end{table}

\begin{table}[h]
\centering
\caption{Projected model rankings from centroid scores.}
\label{tab:model-rankings}
\begin{tabular}{lll}
\toprule
Model & Mean projected score & Responses \\
\midrule
gpt-5.2 & 3.080 & 6 \\
google/gemini-3-flash & 3.077 & 6 \\
gpt-5.4 & 3.027 & 6 \\
claude-sonnet-4-6 & 2.867 & 6 \\
openai/gpt-4.1-mini & 2.853 & 6 \\
deepseek-ai/deepseek-r1 & 2.853 & 6 \\
google/gemini-3-pro & 2.767 & 6 \\
claude-haiku-4-5-20251001 & 2.463 & 6 \\
anthropic/claude-4.5-sonnet & 2.400 & 6 \\
meta/meta-llama-3-70b-instruct & 0.457 & 6 \\
\bottomrule
\end{tabular}

\end{table}

\subsection{Question Perturbation Validation}

The perturbation loop was added as a separate validation lane. Frank variations
are converted into executable question tracks, responses retain their
\texttt{track\_id}, Dasha clusters each track separately, and the perturbation
report compares each variant's dominant reasoning signal against the original
question. Table \ref{tab:perturbation-validation} shows the reported
perturbation smoke: a surface party-name edit preserved the original answer
path, while a material no-writing edit changed the dominant answer path.

\begin{table}[h]
\centering
\caption{Perturbation validation smoke test.}
\label{tab:perturbation-validation}
\begin{tabular}{llll}
\toprule
Track & Type & Comparison & Status \\
\midrule
V\_surface\_invariant & invariant & invariant\_preserved & pass \\
V1 & material & material\_difference\_observed & pass \\
\bottomrule
\end{tabular}

\end{table}

\subsection{Controlled 500-Response Dasha Regression Test}

The 500-response suite is now treated as a deterministic regression test rather
than a robustness demonstration. It generates fixture answers across eight
legally distinct Statute-of-Frauds reasoning archetypes so the software can test
cluster bookkeeping, centroid projection, table generation, and metrics at the
target scale. It does not show that Dasha can discover reasoning families from
natural model outputs.

\begin{table}[h]
\centering
\caption{Controlled Dasha clustering regression test.}
\label{tab:stress-validation}
\begin{tabular}{ll}
\toprule
Check & Value \\
\midrule
Responses & 500 \\
Expected reasoning archetypes & 8 \\
Observed Dasha clusters & 8 \\
Cluster purity & 1.0 \\
Cluster completeness & 1.0 \\
Macro-F1 & 1.0 \\
Mean cluster size & 62.5 \\
\bottomrule
\end{tabular}

\end{table}

The regression test produced eight observed clusters for eight expected reasoning
archetypes. Cluster purity, cluster completeness, and macro-F1 were all 1.0.
This proves that the clustering layer preserves normalized legal-reasoning
distinctions at the target 500-response scale when the signatures are already
known. It does not substitute for the natural-response audit; instead, it
functions as a regression and scale test for cluster bookkeeping, centroid
projection, and metrics generation.

\subsection{Internal Test Coverage}

The reference implementation has 65 research-pipeline tests covering offline run-bundle
exports, Frank packet quality, Karthic rubric quality, SOF gate regression
cases, Dasha clustering invariants, sparse-response gate context, LLM reasoning
signature clustering, source-derived Dasha gate aliasing, a non-SOF
contract-interpretation transfer smoke, judge projection and model ranking,
low-margin Zak escalation, repeated-judge and panel-judge stability aggregation,
Dasha member/centroid audit generation, protocol-freeze manifests, manuscript reference and hygiene linting, model roster config loading, model-call polling, natural-response audit
generation, live-config preflight, credential-alias checks, live-call budget caps, runtime budget enforcement, run-bundle integrity auditing, secret and local-path linting, internal review packet generation, claim-ledger generation, handoff-manifest generation, consolidated no-call audit, method-readiness reporting, review-readiness table generation, metrics, perturbation-track generation, perturbation reporting, and
internal stress-table generation. These tests make the pre-expert-review
manuscript reproducible: the reported claims are backed by executable checks
rather than by visual inspection alone.
